% A fictional one-page resume in the LaTeX subset buildtex.texin imports.
% The person, places and projects are invented for the tests.
\documentclass[11pt]{article}
\usepackage[margin=0.75in]{geometry}
\usepackage{hyperref}
\title{Avery Example}
\author{Portland, Oregon | avery@example.org | example.org \\ github.com/avery-example}
\begin{document}
\maketitle

\section*{Summary}
Software builder who writes \emph{test harnesses} and release tooling, with ten years of field operations
before that. Ships small command-line tools with \textbf{planted-fault checks} and plain-text receipts.

\section{Experience}
\subsection{Release Engineer | Example Tools}
\begin{itemize}
  \item Built a release pipeline in Python that signs every artifact and records its SHA-256; 40 releases
        shipped with no manual step.
  \item Wrote 1,200 regression cases for a parser, each paired with a planted fault that must fail.
  \item Cut median CI time from 14 to 6 minutes by splitting the suite by changed module.
\end{itemize}
\subsection{Field Operations Lead | Riverside Grounds}
\begin{itemize}
  \item Ran crews of four to six on site assessments, estimates and safety briefings, 2014--2024.
  \item Kept the scope a client signed and the work done on site in step, with written change orders.
\end{itemize}

\section{Projects}
\begin{itemize}
  \item \textbf{tally}: a command-line counter for log lines with \texttt{--json} output;
        on PyPI as \texttt{tally-lines}. \href{https://example.org/tally}{example.org/tally}
  \item \textbf{gauge}: renders benchmark intervals as text tables, so a reader sees the spread,
        not only the mean.
\end{itemize}

\section{Skills}
Python, Rust, C, shell; Git and CI; packaging for PyPI and crates.io; technical writing.

\section{Education}
Self-directed study in compilers, typesetting and numerical methods.
\end{document}
